%% kapa: exam=1 round=22 year=2011 period=2 form=A subject=appraisal_law subject_ko=부동산관계법규 seq=19 number=59
%% source: Q-Net 감정평가사 기출문제 (cst00309) artlSeq=5036790 file=extracted/1008/5036790/2011년 제22회 감정평가사 1차시험 기출문제/2교시.hwp
%% conversion: kordoc 4.24.51 -> tools/gen1.py (auto); review=proofread-gpt-astra; flags=none
\begin{question}{59}
부동산 가격공시 및 감정평가에 관한 판례의 입장으로서 옳지 않은 것은?
\begin{choices}
  \item 감정평가서에는 평가원인을 구체적으로 특정하여 명시하여야 하는 동시에 각 요인별 참작 내용과 정도가 객관적으로 납득이 갈 수 있을 정도로 설명되어야 한다.
  \item 수용보상금의 증액을 구하는 소송에서 선행처분으로서 그 수용대상 토지 가격 산정의 기초가 된 비교표준지공시지가결정의 위법을 독립한 사유로 주장할 수 있다.
  \item 보상금의 증감에 관한 소송에 있어서 동일한 사실에 대하여 수개의 감정평가가 있고 그 중 어느 하나의 감정평가가 오류가 있음을 인정할 자료가 없는 이상 법원이 각 감정평가 중 어느 하나만을 채용하는 것은 그 자체로 위법하다.
  \item 개별공시지가는 당해 토지를 담보로 제공함에 있어 그 담보가치를 보장하는 등의 구속력을 가지는 것은 아니다.
  \item 개별공시지가의 결정이 관련 법령이 정하는 절차와 방법에 따라 이루어진 것이라면 그 공시지가가 감정가액이나 실제 거래가격을 초과한다고 하여 위법하게 되는 것은 아니다.
\end{choices}
\end{question}
\answer{3}
